% !TEX root = ../main.tex

% 学士学位论文大摘要
% 相较于前置摘要，大摘要需要更完整地呈现研究背景、动机、方法、实验与结论。
% 本文件在前置中英文摘要之外，补充了对三章核心工作和系统验证结果的展开介绍。

% 2025 版写作规范要求末尾保留英文大摘要，这里保留中文内容但不参与排版。
\iffalse
\begin{digest}[zh]

\subsection*{研究背景与动机}

视觉目标导航是具身智能中连接感知与决策的核心任务。给定机器人当前视觉观察与目标图像，
策略需要在无高精度地图、无全局定位先验的前提下输出可执行的局部运动，
从而完成从起点到目标视角的自主移动。近年来，以 NoMaD 为代表的通用导航模型
将扩散策略引入视觉导航，把未来局部轨迹建模为条件分布，通过去噪采样生成多模态候选，
相比回归单步速度的端到端策略具有显著更强的表达能力。但现有工作的闭环验证主要集中于
轮式平台（如 LoCoBot、TurtleBot、Jackal），面向四足机器人的部署链路仍缺乏系统研究。
四足平台在步态抖动、控制频率、身体姿态耦合等方面与轮式平台差异明显，
这为高层扩散策略与中低层控制接口的协同提出了新的问题。

\subsection*{研究目标与技术路线}

本文以 NoMaD 为高层视觉导航核心，以云深处 Lite3 四足机器人为目标部署平台，
回答三个面向部署的问题：第一，NoMaD 的视觉编码器、距离预测头与扩散策略头在代码
实现上如何耦合，推理层有哪些优化空间；第二，是否可以在不重新训练模型的前提下，
通过采样加速、条件引导、测试时搜索与视觉编码器替换等手段显著提升策略可用性；
第三，这些优化能否整合进面向 Lite3 的分层闭环系统，并完成 MuJoCo 仿真与
Jetson Orin 真机部署的完整验证。技术路线沿``算法层优化—系统层集成—部署层验证''
三级递进展开，对应论文第三、四、五章。

\subsection*{主要工作与创新点}

\textbf{算法层}：建立了统一的离线统计实验链路，对 DDIM 推理加速、Classifier-Free
Guidance (CFG) 引导增强、Test-Time Scaling (TTS) 推理时搜索与视觉编码器升级进行了
系统对比。DDIM 实验中 \texttt{ddim\_2} 配置相对 DDPM-10 基线实现约 4.84 倍加速且
轨迹质量无显著退化；CFG 实验通过 score function 修改分析给出引导尺度选择准则，
指出其有效范围有限、宜作为可调参数；TTS 实验提出将总收益分解为``评分筛选''
（近乎零时延代价）与``候选池扩容''（时延线性增长但边际递减）两部分，TTS-8 在三因素
联合排序中以最小代价取得综合最优，为推理预算分配提供了新的诠释；视觉编码器升级
采用两阶段框架比较 EfficientNet-B0、ResNet-50、DINOv2-Small、ConvNeXt-Tiny 的
离线指标，为后续部署筛选出候选组合。

\textbf{系统层}：将统一推理模块、状态机、导航主机、中层控制映射与平台接口整合为
分层闭环系统，支持 MuJoCo 仿真与真机两种后端。状态机由 idle / stand / navigate /
explore / recovery / estop / completed / failed 组成，通过统一接口与平台解耦；
导航主机支持批量计划（JSON）和交互式任务两种模式，允许在同一会话内完成拍照、建图、
导航、探索等组合任务。MuJoCo 闭环实验中，Lite3 在 easy/medium/hard 三张地图上均
成功完成定点导航、多目标导航、自由探索任务，所有目标均被到达。

\textbf{部署层}：针对 Jetson Orin 与 Lite3 上位机链路，形成了分阶段联调流程、
多级安全机制与图像记录规范。Orin 阶段一（独立调试）验证了相机采集、模型加载、
推理基准与端到端流水线；在桥接层对速度命令加入 NaN/Inf 守卫、幅值限幅与
连接健康监测，任意推理异常自动走安全急停；状态机层加入全局摔倒保护，使
idle/stand 等状态也不再遗漏摔倒判定。

\subsection*{实验与结论}

离线实验表明，DDIM-2 + 可选 CFG + 可配置 TTS 是当前最值得进入闭环验证的组合；
MuJoCo 闭环表明分层系统能够稳定完成四足导航与探索任务，多目标任务中的任务切换
不会破坏视觉上下文；Orin 阶段一的独立测试验证了真机流水线的基础可用性。
本文将 NoMaD 从轮式平台的实验框架扩展为面向四足平台的可分析、可仿真、可部署的
研究链路，并沉淀出一套可复用的推理优化实验协议、闭环系统组织方式以及真机部署
安全护栏清单，为后续基于扩散策略的四足机器人导航研究提供了参考基础。

\end{digest}
\fi

\begin{digest}[en]

\subsection*{Background and Motivation}

Visual goal navigation is a core task in embodied intelligence that connects perception
with decision-making. Given the current visual observation and a goal image, the policy
must output executable local motion without relying on high-precision maps or global
localization, thereby achieving autonomous movement from the current viewpoint to the
goal viewpoint. As robots are increasingly deployed in unstructured open-world
environments such as campus patrol, indoor guidance, last-mile delivery and home
service, navigation can no longer be reduced to a pure geometric path-planning problem
on a known map. It must instead handle weak priors, partial observability, dynamic
obstacles and ambiguous goal specifications, which has motivated a paradigm shift from
classical SLAM-plus-planner pipelines toward learning-based, image-conditioned policies
that consume raw visual observations directly.

Recently, general navigation models such as GNM, ViNT, and NoMaD have introduced
diffusion policies into visual navigation, modelling the future local trajectory as a
conditional distribution and generating multi-modal candidates through iterative
denoising sampling. This formulation has significantly stronger expressive power than
regression-based end-to-end policies that collapse multi-modal decisions (such as
turning left or right at a T-junction) into a single mean output. By modelling the
full distribution of feasible local trajectories, diffusion-based policies are able to
preserve multi-modality, sample diverse candidates, and integrate naturally with
test-time selection and guidance mechanisms. However, closed-loop validation of these
models has so far concentrated on wheeled platforms (LoCoBot, TurtleBot, Jackal),
while deployment pipelines targeting quadruped robots remain underexplored. Quadruped
platforms differ substantially from wheeled ones in gait-induced jitter, control-loop
frequency, body-posture coupling, foot-impact disturbance, and real-time safety
requirements. These differences pose new challenges for the collaboration between
high-level diffusion policies and mid-/low-level control interfaces, and they cannot
be addressed by simply re-routing the wheeled-robot velocity command to a quadruped's
locomotion controller.

\subsection*{Objectives and Approach}

This thesis takes NoMaD as the high-level navigation core and the DeepRobotics Lite3
quadruped as the target deployment platform, and addresses three deployment-oriented
questions: (i) how the visual encoder, distance-prediction head, and diffusion policy
head are coupled in the NoMaD implementation and where inference-level optimizations
are possible without retraining the model; (ii) whether, without retraining, sampling
acceleration via DDIM, conditional guidance via Classifier-Free Guidance (CFG),
test-time search via Test-Time Scaling (TTS), and visual encoder replacement can
meaningfully improve policy usability for resource-constrained edge deployment; and
(iii) whether these optimizations can be integrated into a hierarchical closed-loop
system for Lite3 and validated through MuJoCo simulation as well as real-world
deployment on the Jetson Orin onboard computer. The technical roadmap proceeds through
algorithm-level optimization, system-level integration, and deployment-level
validation, corresponding respectively to Chapters 3, 4, and 5 of the thesis. A
deliberate emphasis is placed on bridging the offline-evaluation versus closed-loop
deployment gap, ensuring that every algorithmic conclusion is re-examined under
realistic timing, safety, and platform constraints.

\subsection*{Main Contributions}

\textbf{Algorithm layer.} A unified offline statistical pipeline is built to compare
DDIM inference acceleration, Classifier-Free Guidance (CFG) enhancement, Test-Time
Scaling (TTS) search, and visual encoder upgrades under a single evaluation protocol.
Concretely, DDIM-2 attains about $4.84\times$ speedup over the DDPM-10 baseline on the
GoStanford evaluation set without notable trajectory-quality degradation, with the
mean-squared deviation against the baseline trajectory distribution remaining below
$0.022$. The CFG analysis, derived from score-function modification, establishes a
guidance-scale selection rule and shows that CFG roughly doubles the per-step compute
while delivering only marginal Forward-Progress gain in the diffusion-policy regime;
hence CFG is best treated as a tunable knob rather than a default-on component. The
TTS analysis decomposes total gain into a nearly zero-latency \emph{scoring-based
selection} component (which re-ranks the eight candidates already produced by a single
diffusion batch) and a linearly priced but rapidly diminishing \emph{candidate-pool
expansion} component (which spawns extra batches to enlarge the pool). TTS-8 wins the
three-factor joint ranking at minimal cost, providing a novel interpretation of
inference-budget allocation as ``re-rank cheaply first, expand only when necessary.''
Encoder upgrades use a two-stage framework to compare EfficientNet-B0, ResNet-50,
DINOv2-Small, and ConvNeXt-Tiny on identical offline cases; the encoder $\times$ DDIM
$/$ CFG $/$ TTS joint ablation establishes that DDIM-2 and CFG default-off remain
robust across all four encoders, while TTS gain scales with the initial candidate
diversity of the chosen encoder.

\textbf{System layer.} A hierarchical closed-loop system---comprising a unified
inference module, an explicit finite-state machine, a navigation host, a mid-level
waypoint-to-velocity mapping, and a platform interface---supports both MuJoCo and
real-robot backends through one set of high-level code. The state machine
(\texttt{idle / stand / navigate / explore / recovery / estop / completed / failed})
is decoupled from the physical platform via a uniform interface, so that switching
between MuJoCo and Lite3 only changes the backend implementation rather than the
high-level logic. The navigation host supports both batched JSON plans and interactive
single commands, allowing capture, topomap building, navigation, and exploration to be
composed within a single experiment session. Mid-level mapping converts the diffusion
policy's local waypoints into bounded forward, lateral, and yaw velocity commands,
respecting Lite3-specific safety envelopes ($v_x \le 0.4\,\text{m/s}$ in simulation,
$v_x \le 0.12\,\text{m/s}$ during real-robot deployment). In MuJoCo closed-loop
experiments, Lite3 successfully completes point navigation, multi-goal navigation, and
free exploration on easy / medium / hard maps, reaching every goal across repeated
runs. The same unified inference module that handles the four offline optimizations is
re-used at deployment time, demonstrating that the algorithm-layer findings translate
directly into operational system behaviour.

\textbf{Deployment layer.} For the Jetson Orin onboard host and the Lite3
upper-computer link, the thesis develops a staged commissioning procedure, multi-level
safety mechanisms, and an image-recording convention. Stage 1 (Orin standalone
commissioning) validates camera capture, model loading, inference benchmarking, and
the end-to-end pipeline in isolation from the robot, so that any inference-side defect
is discovered before the platform is energized. Stage 2 (Orin and Lite3 joint
commissioning) progresses from network and heartbeat verification, through standup and
low-speed twist tests, to interactive and goal-directed navigation under conservative
speed limits. The bridging layer wraps the Lite3 SDK with a 25\,Hz velocity-refresh
thread and a 5\,Hz heartbeat thread, repeating the most recent clipped velocity
between higher-level NoMaD inference cycles to guarantee continuous control flow even
when the policy runs at single-digit hertz. Multi-level safety guards include
NaN/Inf clipping, magnitude limits, fallen-robot detection, and a connection-stale
threshold of two seconds; any inference anomaly or platform fault is routed to a soft
emergency stop. Real-robot experiments on multiple indoor scenes with pre-collected
topomaps record closed-loop frequencies of 2.89\,Hz (DDPM baseline), 6.80\,Hz
(DDIM-2 + CFG=0 + TTS-8), 3.43\,Hz (DDIM-2 + CFG=2 + TTS-8), and 5.90\,Hz
(DDIM-3 + CFG=0, no TTS), with a 100\% success rate across all four configurations
under the present conservative speed setting. The frequency ranking is fully
consistent with the offline analysis in Chapter 3: DDIM step count dominates latency,
CFG approximately halves the achievable rate, and TTS-8's overhead is negligible
because re-ranking eight already-generated candidates does not require an extra
diffusion batch.

\subsection*{Experiments and Conclusions}

Offline results indicate that DDIM-2 plus optional CFG and configurable TTS-8 is the
most promising configuration to bring into closed-loop validation, and the encoder
upgrade pipeline supplies a reusable framework for future visual-backbone exploration
rather than a single best-encoder verdict. MuJoCo experiments confirm that the
hierarchical system can robustly complete quadruped navigation and exploration tasks,
that multi-goal task switching preserves visual context, and that the recovery
sub-state successfully unblocks the robot from corridor-corner stalls without
requiring task abortion. Stage-1 validation on Jetson Orin demonstrates the baseline
feasibility of the real-robot pipeline (about 7\,Hz benchmarked NoMaD inference,
$\approx 151$--$157\,\mathrm{ms}$ steady-state pipeline cycle once warm-up cost is
amortized), and Stage-2 multi-scene results provide the first real-robot evidence
that diffusion-policy-based image-goal navigation can run end-to-end on a Lite3
quadruped at single-digit hertz while still completing every commissioned trial.
Crucially, the fact that even the slow 2.89\,Hz baseline achieves 100\% success
demonstrates that the closed-loop system has substantial robustness margin: the
conservative 0.12\,m/s speed cap means each cycle's displacement is only about
4\,cm---far smaller than the topomap node spacing---so high-level latency tolerance is
absorbed by the topomap sliding window and the 25\,Hz bridge replay rather than
manifesting as catastrophic drift.

\subsection*{Limitations and Future Work}

Three limitations are explicitly acknowledged. First, the encoder upgrade conclusions
remain offline statistics plus simulator validation, and a final
real-robot-conditioned encoder ranking will require benchmarking inference latency,
image-distribution robustness, and failure modes directly on the Orin/Lite3 stack with
real cameras. Second, the present real-robot statistics, while sufficient to show that
all four inference configurations succeed under the current scenario scale, do not yet
include large-sample success-rate confidence intervals or formal trajectory-deviation
distributions; expanding the experiment matrix to outdoor scenes, dynamic obstacles,
and longer corridors should enhance differentiation between configurations and unlock
the latency margin that the fastest scheme has reserved. Third, the topomap is still
collected offline by an operator; an online maintenance loop that refreshes nodes when
the environment changes would improve long-term autonomy. Future work will therefore
focus on (i) finalizing the on-robot encoder choice with end-to-end real-world
metrics, (ii) running large-scale repeated trials with formal statistics, (iii)
generalizing the unified inference module and state machine to additional platforms
such as the Tron1 humanoid robot to test cross-embodiment portability, and (iv)
extending image-goal conditioning to language-grounded targets, allowing operators to
issue natural-language commands such as ``go to the red door at the end of the
corridor'' that get translated into NoMaD-compatible visual conditions.

The thesis extends NoMaD from a wheeled-platform research framework into an
analyzable, simulatable, and deployable research pipeline for quadruped platforms,
distilling a reusable inference-optimization protocol, a hierarchical closed-loop
system organization, and a real-robot safety-guard checklist. Together, these
artefacts provide a practical basis for subsequent research on diffusion-policy-based
visual navigation for quadruped robots, and they lower the engineering threshold for
follow-up work that wishes to validate algorithmic ideas under genuine on-robot
conditions rather than only on offline benchmarks.

\end{digest}
